% ===========================================================================
%  7 — Pipeline de dados
% ===========================================================================
\chapter{Pipeline de dados e canonicalização}
\label{cap:pipeline}

\section{Caracterização da fonte de dados}

O conjunto de dados original é um arquivo de \SI{31}{\mega\byte} com telemetria
de uma bancada instrumentada de máquina rotativa.

\begin{table}[H]
\caption{Caracterização do conjunto de dados \arq{banner.csv}}
\label{tab:dataset}
\begin{tabularx}{\linewidth}{@{}L{46mm}X@{}}
\toprule
\headrow \thd{Propriedade} & \thd{Valor}\\
\midrule
Registros & \num{166796}\\
\zebra Período & 30/04/2026 a 16/06/2026 (aproximadamente 47 dias)\\
Colunas de sensor utilizadas & 18 (em unidades SI)\\
\zebra Colunas descartadas & Duplicatas em \si{\inch\per\second} e
  \si{\degF} --- colineares por construção com as contrapartes métricas\\
Rótulos brutos distintos & \num{151}\\
\zebra Famílias canônicas & 17 (12 falhas + 5 estados operacionais)\\
Campanhas com $\geq 30$ eventos & 146\\
\zebra Regimes de rotação &
  \SI{500}{rpm} (\num{55857}) $\cdot$
  \SI{2000}{rpm} (\num{55160}) $\cdot$
  \SI{1000}{rpm} (\num{53414}) $\cdot$
  \SI{3000}{rpm} (\num{1707}) $\cdot$
  \SI{0}{rpm} (658)\\
Equipamentos & 1 (bancada única)\\
\bottomrule
\end{tabularx}
\end{table}

\begin{limitacao}
Os regimes de rotação são fortemente desbalanceados: três regimes concentram
\SI{98.6}{\percent} dos registros, e o regime de \SI{3000}{rpm} --- justamente
onde o desbalanceamento se manifesta de forma inequívoca
(Seção~\ref{sec:fisica}) --- tem apenas \num{1707} amostras. Toda conclusão
sobre alta rotação repousa sobre \SI{1}{\percent} dos dados.
\end{limitacao}

\section{Canonicalização de rótulos}

\subsection{O problema}

A coluna de anotação do conjunto de dados contém \num{151} valores distintos
para 17 condições reais. A dispersão tem três causas, todas presentes em dados
industriais reais:

\begin{enumerate}
  \item \textbf{Variação de campanha}: sufixos e prefixos que identificam a
    sessão de coleta, não a falha --- \cd{\_2}, \cd{\_3}, \cd{\_pos\_2},
    \cd{\_carga}, \cd{\_adxl\_0}, \cd{new\_}.
  \item \textbf{Erros de digitação}: \cd{ddesbalanceado},
    \cd{cockecocked\_adxl\_0}, \cd{mortor\_desligado}, \cd{normla\_carga\_3\_3},
    \cd{desbanlanceado}, \cd{desabalanceado}, \cd{desabanceado}.
  \item \textbf{Mistura de falhas com estados}: \cd{baseline}, \cd{teste},
    \cd{acelerando} e \cd{motor\_desligado} não são defeitos, e tratá-los como
    tal produziria prescrições para condições normais.
\end{enumerate}

\subsection{A solução: regras ordenadas e auditáveis}

A canonicalização é feita por uma lista \emph{ordenada} de expressões
regulares declaradas em YAML. A primeira regra que casa vence. Regras de falha
precedem regras de estado, para que \cd{rolamento\_outer\_novo\_teste} caia em
\cd{rolamento\_outer} e não em \cd{teste}.

\begin{lstlisting}[language=yamlpmx, caption={Extrato de \arq{src/etl/label\_map.yaml} --- declaração de famílias e regras ordenadas}, label={lst:labelmap}]
families:
  rolamento_inner:  { is_fault: true,  display: "Defeito na pista interna do rolamento (BPFI)" }
  desbalanceamento: { is_fault: true,  display: "Desbalanceamento do rotor" }
  cocked_rotor:     { is_fault: true,  display: "Rotor inclinado (cocked rotor)" }
  normal:           { is_fault: false, display: "Operacao normal" }
  motor_desligado:  { is_fault: false, display: "Motor desligado" }

# A ordem importa: a primeira regra que casar vence.
rules:
  - { pattern: "inner",                  family: rolamento_inner }
  - { pattern: "outer",                  family: rolamento_outer }
  - { pattern: "ball",                   family: rolamento_ball }
  - { pattern: "comb",                   family: rolamento_combination }
  # cobre desb-, e os erros reais ddesb-, desbanlanceado, desabalanceado
  - { pattern: "desb|abal|abanc|balanc", family: desbalanceamento }
  - { pattern: "desalinh",               family: desalinhamento }
  - { pattern: "cock",                   family: cocked_rotor }   # cobre "cockecocked"
  # --- estados operacionais, avaliados por ultimo ---
  - { pattern: "desligado",              family: motor_desligado } # cobre "mortor_"
  - { pattern: "norm",                   family: normal }          # cobre "normla"
  - { pattern: "^(new_)?tes",            family: teste }
\end{lstlisting}

\begin{decisao}
A alternativa óbvia --- agrupamento automático por similaridade de texto ---
foi rejeitada por duas razões. Primeiro, não é auditável: ninguém consegue
justificar, meses depois, por que dois rótulos caíram no mesmo grupo.
Segundo, erro de grafia não é problema de similaridade semântica:
\cd{cockecocked} e \cd{cocked} são vizinhos por acidente de digitação, e
\cd{rolamento\_inner} e \cd{rolamento\_outer} são textualmente próximos mas
descrevem defeitos distintos, com procedimentos distintos. Regras explícitas
erram menos e explicam-se sozinhas (\adrr{ADR-05}).
\end{decisao}

\subsection{Auditoria: falhar cedo em vez de classificar em silêncio}

O canonizador levanta exceção diante de rótulo sem regra. O ETL, por sua vez,
canonicaliza \textbf{todos} os rótulos distintos \emph{antes} de gravar
qualquer linha:

\begin{lstlisting}[language=Python, caption={Auditoria integral antes da gravação --- \arq{scripts/ingest\_data.py}}, label={lst:auditoria}]
# Canonicalizacao: audita 100% dos rotulos antes de gravar qualquer coisa
raw_labels = df["fault"].astype(str).str.strip().str.lower()
unique_labels = sorted(raw_labels.unique())
mapping = {lbl: canonizer.canonize(lbl) for lbl in unique_labels}  # levanta se desconhecido
print(f"{len(unique_labels)} rotulos brutos mapeados para "
      f"{len({f.name for f in mapping.values()})} familias canonicas.")
\end{lstlisting}

A consequência prática: um rótulo novo em uma carga futura interrompe o
processo com o nome exato do rótulo, exigindo curadoria humana. É o oposto do
comportamento usual --- descartar a linha ou atribuí-la a uma categoria
\enquote{outros} --- que produz uma base silenciosamente incorreta
(\req{RNF-14}).

\section{Fluxo do ETL}

\begin{figure}[H]
\centering
\fitwidth{%
\begin{tikzpicture}[x=1cm, y=1cm]

  \node[etapaio, text width=27mm] (csv) at (1.9,0)
    {\arq{banner.csv}\\\SI{31}{\mega\byte}\\\num{166796} linhas};
  \node[etapa, text width=28mm] (rd) at (6.0,0)
    {\textbf{Leitura e tipagem}\\{\scriptsize \cd{parse\_dates}, coerção
     numérica}};
  \node[etapa, text width=28mm] (cn) at (10.2,0)
    {\textbf{Canonicalização integral}\\{\scriptsize \num{151} rótulos
     $\rightarrow$ 17 famílias}};
  \node[etapa, text width=28mm] (dd) at (14.4,0)
    {\textbf{Deduplicação}\\{\scriptsize por \cd{id}; descarte de linhas
     incompletas}};

  \node[etapa, text width=28mm] (tx) at (14.4,-3.6)
    {\textbf{Taxonomia}\\{\scriptsize popula \cd{fault\_families} e
     \cd{label\_map}}};
  \node[etapa, text width=28mm] (ld) at (10.2,-3.6)
    {\textbf{Carga dos eventos}\\{\scriptsize inserção idempotente em
     \cd{events}}};
  \node[etapaout, text width=27mm] (pg) at (6.0,-3.6)
    {\textbf{PostgreSQL}\\{\scriptsize histórico consultável}};
  \node[etapaerr, text width=27mm, fill=pmxOrangeWash, draw=pmxOrange]
       (err) at (1.9,-3.6)
    {\textbf{Interrupção}\\{\scriptsize rótulo sem regra $\rightarrow$
     exceção nomeando o rótulo}};

  \draw[pipe] (csv) -- (rd);
  \draw[pipe] (rd) -- (cn);
  \draw[pipe] (cn) -- (dd);
  \draw[pipe] (dd) -- (tx);
  \draw[pipe] (tx) -- (ld);
  \draw[pipe] (ld) -- (pg);
  \draw[pipe, draw=pmxOrange] (cn.south) -- (10.2,-2.2) -- (1.9,-2.2) -- (err.north)
    node[pipelbl, pos=0.55, above=1pt]{rótulo desconhecido};

\end{tikzpicture}}
\caption[Pipeline do ETL]{Pipeline do ETL do histórico. A canonicalização
  integral acontece \emph{antes} de qualquer escrita: se algum rótulo não casar
  com regra, o processo termina sem tocar o banco. O caminho em laranja é a
  porta de qualidade de dados.}
\label{fig:etl}
\end{figure}

\section{\textit{Features} geradas na carga}

O ETL grava as 18 grandezas brutas; as 5 \textit{features} derivadas são
calculadas em tempo de treino e de inferência pelo \emph{mesmo} módulo
(\arq{src/ml/features.py}), nunca materializadas no banco.

\begin{decisao}
Não materializar as derivadas evita a classe de defeito mais insidiosa em
sistemas de aprendizado de máquina: a divergência entre a transformação aplicada
no treino e a aplicada na inferência --- o \textit{training-serving skew}
descrito por Sculley \textit{et al.}~\cite{sculley2015debt}. Com uma única
função compartilhada, a divergência é impossível por construção. O custo é
recalcular cinco divisões por evento, o que é irrelevante.
\end{decisao}

\section{Agregações analíticas}

As análises do painel são consultas SQL executadas sob demanda, não artefatos
pré-computados. Três agregações não triviais merecem registro:

\begin{tabularx}{\linewidth}{@{}L{34mm}X@{}}
\toprule
\headrow \thd{Agregação} & \thd{O que calcula e por quê}\\
\midrule
\cd{severity\_by\_rpm} & Percentis 50, 95 e 99 da velocidade eficaz por família
  \emph{e} por regime de rotação, mais a linha de base do estado normal no mesmo
  regime. Usa percentis em vez de média porque as distribuições têm cauda longa
  e a média seria puxada por transientes de partida.\\
\zebra \cd{fault\_signatures} & Escore padronizado de cada família em cada
  \textit{feature} e a razão entre variância inter e intrafamília --- no espírito
  da estatística $F$ da análise de variância. Mostra objetivamente quais
  indicadores separam as falhas.\\
\cd{leakage\_evidence} & Decomposição da dispersão de cada \textit{feature}
  \emph{entre} campanhas de coleta e \emph{dentro} de cada campanha. É a
  evidência quantitativa do confundimento analisado na
  Seção~\ref{sec:vazamento}.\\
\bottomrule
\end{tabularx}

\begin{nota}
As três agregações rodam em torno de \SI{0.2}{\second} sobre \num{166796}
eventos. Por isso não há pré-cálculo: o painel reflete inclusive os eventos que
chegam por MQTT durante uma demonstração ao vivo. Interpolação de nomes de
coluna nas consultas é protegida por lista branca derivada da constante de
\textit{features} do próprio código --- coluna fora da lista levanta exceção.
\end{nota}
